\chapter{Estado del arte}
\label{cap:estado-arte}

\section{La diabetes mellitus tipo 1 y el monitoreo continuo de glucosa}

La diabetes mellitus tipo 1 es una patología crónica en la que el organismo deja de
producir insulina, lo que obliga al paciente a tomar decisiones terapéuticas de forma
continua. En el día a día esto se traduce en medir la glucosa, calcular dosis de
insulina, estimar hidratos de carbono y anticipar el efecto del ejercicio. El problema no
es únicamente clínico, sino también informacional: el usuario debe interpretar datos muy
heterogéneos en un intervalo temporal corto.

Los sistemas de monitorización continua de glucosa suponen un salto cualitativo frente a
la glucemia capilar tradicional, ya que permiten disponer de un valor reciente de glucosa
intersticial y, además, de una tendencia de evolución. Esa combinación de valor actual,
histórico y dirección de cambio hace posible detectar hipoglucemias, hiperglucemias o
subidas rápidas antes de que el problema sea evidente desde una única medición aislada.

Sin embargo, disponer de más datos no resuelve por sí mismo la dificultad de gestión. El
usuario sigue necesitando relacionar esas lecturas con la insulina administrada, la comida
ingerida y la actividad física realizada. Precisamente ahí aparece el hueco que pretende
cubrir GlucoMate: no solo leer glucosa, sino contextualizarla.

\section{Sistemas CGM comerciales: Abbott Freestyle Libre}

Dentro de los sistemas comerciales de CGM (\textbf{monitoreo continuo de glucosa}), la familia \textit{Freestyle Libre} se ha
consolidado como una de las más extendidas por su facilidad de uso y por la integración
con teléfonos móviles. El sensor se adhiere al brazo y permite recuperar datos de glucosa
mediante tecnologías inalámbricas, principalmente NFC para ciertas interacciones de
lectura y Bluetooth Low Energy para la comunicación continuada en modelos recientes.

Para este proyecto, el sensor de referencia es el \textit{Freestyle Libre 2 Plus}. La
elección no es casual: es un dispositivo ampliamente utilizado, con una comunidad activa
de usuarios y desarrolladores interesados en comprender su funcionamiento más allá del
ecosistema oficial. Además, representa un caso realista de integración entre hardware
médico, protocolos propietarios y aplicaciones móviles de terceros.

La principal limitación del entorno comercial no es la capacidad de medición, sino el
grado de apertura. Las aplicaciones oficiales priorizan la seguridad del ecosistema y la
experiencia cerrada del fabricante, pero ofrecen poca flexibilidad para integrar la señal
glucémica con otros datos generados por el paciente o para experimentar con nuevas capas
de análisis. GlucoMate se plantea precisamente como una alternativa investigadora sobre
esa base tecnológica.

\section{Aplicaciones existentes: xDrip+, LibreMonitor, Glimp}

Antes de plantear una solución propia es necesario analizar las herramientas ya
existentes. Proyectos como \textit{xDrip+}, \textit{LibreMonitor} o \textit{Glimp} han
demostrado que existe un interés real en ampliar las posibilidades de los sensores Abbott
mediante software no oficial. Estas aplicaciones destacan por ofrecer funciones de
visualización avanzada, alarmas, exportación de datos o integración con otros servicios.

No obstante, muchas de ellas están orientadas a perfiles de usuario avanzados, requieren
configuraciones técnicas poco amigables o se centran en un subconjunto del problema.
Algunas sobresalen en conectividad, otras en comunidad, y otras en soporte histórico, pero
no todas ofrecen un equilibrio adecuado entre simplicidad de uso, registro clínico,
persistencia estructurada e interpretación pedagógica de los datos.

La oportunidad de GlucoMate no consiste en replicar exactamente esas herramientas, sino
en combinar varias ideas valiosas en una propuesta cohesionada: lectura de glucosa,
registro de comidas, insulina y ejercicio, perfil terapéutico configurable y un agente de
apoyo basado en IA capaz de usar ese contexto de manera integrada.

\section{Ingeniería inversa del protocolo Abbott}

Uno de los aspectos más relevantes del estado del arte en este proyecto es el trabajo de
ingeniería inversa realizado por la comunidad sobre los sensores de Abbott. Aunque el
transporte físico puede apoyarse en estándares conocidos, como ISO~15693 en el ámbito NFC
o BLE en la comunicación continua, la semántica de determinados comandos y el formato de
los datos no forman parte de una documentación pública orientada a desarrolladores de
terceros.

Como parte del análisis experimental se realizó un intento de lectura directa del sensor
\textit{Freestyle Libre 2 Plus} desde un dispositivo Android. Desde el punto de vista de
hardware, el sensor expone una interfaz compatible con NFC-V, tecnología asociada al
estándar ISO~15693. El objetivo de esta prueba era validar si un teléfono móvil podía
establecer comunicación con el parche, identificarlo correctamente y acceder a sus
bloques de memoria mediante comandos NFC estándar.

Para esta primera aproximación se utilizó React Native junto con Expo. Sin embargo, se
comprobó que la aplicación estándar Expo Go no resultaba suficiente para este caso de
uso, ya que no incorpora los módulos nativos necesarios para acceder a la antena NFC del
dispositivo Android. Por este motivo, fue necesario plantear una \textit{Development
Build} mediante EAS Build, incorporando la librería \texttt{react-native-nfc-manager} y
los permisos específicos de Android requeridos para el uso de NFC.

Durante la puesta en marcha del entorno de pruebas apareció una dificultad adicional
relacionada con la conexión entre la aplicación instalada en el móvil y el servidor de
desarrollo Metro. El APK intentaba resolver \texttt{localhost} o \texttt{127.0.0.1} desde
el propio dispositivo Android, lo que impedía alcanzar el servidor ejecutado en el
ordenador. A esto se sumaban posibles restricciones de red, como el cortafuegos de
Windows o la falta de visibilidad entre direcciones IP de la red local. Para resolverlo
se utilizó ADB como puente de depuración mediante USB. Tras autorizar el dispositivo, se
redirigió el puerto de Metro con el comando \texttt{adb reverse tcp:8081 tcp:8081} y se
lanzó Expo en modo local. De esta forma, las peticiones realizadas por la aplicación a
\texttt{localhost:8081} quedaban encaminadas hacia el servidor Metro del ordenador.

Una vez estabilizado el entorno de ejecución, se implementó una primera prueba de
comunicación NFC con el sensor. El código básico consistió en solicitar explícitamente la
tecnología NFC-V y recuperar la información de la etiqueta detectada:

\begin{lstlisting}[language=JavaScript, caption={Prueba inicial de comunicación NFC-V con el sensor}, label={cod:nfcv-inicial}, captionpos=b]
await NfcManager.requestTechnology(NfcTech.NfcV);
const tag = await NfcManager.getTag();
\end{lstlisting}

La llamada a \texttt{NfcTech.NfcV} indica al sistema Android que la comunicación debe
realizarse utilizando el protocolo ISO~15693. Esta selección es esencial, ya que el
sensor no responde si se intenta acceder a él mediante una tecnología NFC distinta. La
posterior llamada a \texttt{getTag()} permitió obtener información básica del dispositivo,
incluyendo las tecnologías soportadas y su identificador único. En la prueba realizada se
detectó una etiqueta con soporte \texttt{android.nfc.tech.NfcV} y
\texttt{android.nfc.tech.NdefFormatable}, además de un identificador de sensor similar a
\texttt{7328D38A78007AE0}. Este resultado confirmó que el teléfono era capaz de detectar
el parche, reconocerlo como dispositivo NFC-V y establecer una comunicación inicial con
él.

Tras validar el reconocimiento del sensor, se intentó leer su memoria interna mediante el
comando estándar \textit{Read Single Block} de ISO~15693, identificado por el código
\texttt{0x20}. Se realizaron lecturas sucesivas de los bloques comprendidos entre el
bloque 0 y el bloque 42. Aunque el sensor devolvió datos para distintos bloques, estos se
presentaban como secuencias hexadecimales sin una estructura interpretable de forma
directa. Por ejemplo, algunos bloques comenzaban con valores similares a
\texttt{0da19cf20a51f2b7}, lo que evidenciaba que la lectura física era posible, pero no
que el contenido pudiera transformarse directamente en una medida de glucosa.

Este comportamiento coincide con la evolución de las medidas de protección introducidas
por Abbott a partir de generaciones posteriores al \textit{Freestyle Libre} original. En
modelos como el \textit{Freestyle Libre 2 Plus}, los datos almacenados y transmitidos no
pueden interpretarse aplicando únicamente una fórmula sobre la memoria leída en bruto. La
información relevante se encuentra protegida mediante mecanismos criptográficos
propietarios, habitualmente descritos por la comunidad como esquemas basados en cifrado
simétrico y claves derivadas del identificador del sensor y de procesos internos de
emparejamiento. En consecuencia, el problema técnico deja de ser únicamente una cuestión
de acceso NFC y pasa a incluir la comprensión del protocolo propietario, la gestión de
claves y el descifrado de los datos.

La conclusión principal de esta prueba fue que la comunicación NFC-V con el sensor es
técnicamente viable desde Android, pero que la lectura directa de bloques ISO~15693 no es
suficiente para obtener valores de glucosa utilizables. El obstáculo principal no se
encuentra en la detección del sensor ni en el acceso físico a la memoria, sino en la
interpretación criptográfica de la información obtenida. Por ello, el enfoque del
proyecto se orientó hacia el estudio de implementaciones comunitarias ya existentes,
especialmente Juggluco, que actúan como referencia práctica para comprender el proceso de
lectura, emparejamiento e interpretación de sensores Abbott recientes.

En este contexto aparecen referencias frecuentes a comandos propietarios como
\texttt{0xA0} y \texttt{0xA1}, empleados en discusiones técnicas y proyectos de la
comunidad para describir operaciones de acceso a información específica del sensor. En la
práctica, estos comandos no deben entenderse como parte del estándar ISO~15693, sino como
extensiones del fabricante sobre una base de comunicación conocida. Su presencia explica
por qué la simple disponibilidad de NFC en el teléfono no basta para implementar un lector
completo sin conocimiento adicional del protocolo.

La existencia de aplicaciones puente como Juggluco encaja en este escenario. En el
proyecto analizado, Juggluco actúa como intermediario local para obtener la lectura ya
interpretada y ponerla a disposición de la aplicación principal mediante HTTP en la misma
red o en el propio dispositivo. Esta aproximación reduce la complejidad de la lectura
directa y permite centrar el esfuerzo del TFG en la integración software, la persistencia
de datos y la capa de ayuda inteligente.

\section{Normativa y consideraciones legales}

El uso de sensores médicos desde aplicaciones no oficiales introduce cuestiones legales y
éticas que no pueden ignorarse. GlucoMate no es una aplicación oficial del fabricante ni
un producto sanitario certificado, por lo que sus resultados deben entenderse como apoyo
informativo y no como sustitución de la supervisión clínica ni de las herramientas
reguladas.

Además, el proyecto trata datos de salud de especial sensibilidad: lecturas de glucosa,
tratamiento farmacológico, hábitos alimentarios y actividad física. Por ello, aunque se
trate de un prototipo académico, la arquitectura debe respetar principios básicos de
minimización de datos, separación de credenciales, uso controlado de servicios externos y
transparencia respecto al propósito de cada tratamiento.

Desde la perspectiva del TFG, este marco legal no invalida el desarrollo experimental,
pero sí obliga a delimitar claramente su alcance. La aplicación se presenta como una
herramienta de investigación y apoyo, y el agente de IA incorpora avisos explícitos para
recordar que sus recomendaciones son orientativas.